\section{The complex problem}\label{sec:complex}

Over $\C$ the problem is solved, and the answer stops exactly where the
elementary methods of \S\ref{sec:first}--\S\ref{sec:chart} stop:
$\GtzC(m,k)$ holds for $k\ge2$ if and only if $m\le k+1$.  Half of that
is already proved.  Nothing in \S\ref{sec:first} or \S\ref{sec:chart}
consulted the field, so their positive results stand over $\C$
unchanged, and they cover $m\le k+1$ (\S\ref{sec:cpositive}).

The other half needs one counterexample and a way to move it.  A single
design of size four and rank two fails, and two operations propagate the
failure across the board: one raises the size and the rank together, the
other raises the size alone, and each keeps a bound on the least
eigenvalue that any $k$-subset can achieve.

\subsection{The value of a design}\label{sec:cvalue}

Once the answer is \emph{no}, a number remains, and the operations of
\S\ref{sec:cpadding} propagate it.  Fix a
design $\dsn$ at $(m,k)$ over $\FF$.  For every $k$-subset $C$ and every
$x\in\FF^k$,
\begin{equation}\label{eq:cover}
  x\adj\SC x\;=\;\sum_{c\in C}\lvert\langle g_c,x\rangle\rvert^{2},
  \qquad\text{so}\qquad
  \SC\psd s\,I_k
  \iff
  s\lVert x\rVert^{2}\le\sum_{c\in C}\lvert\langle g_c,x\rangle\rvert^{2}
  \ \ \text{for all }x .
\end{equation}
Both readings are used below: the left one to recognize a design, the
right one to transport a bound along a map of index sets.  Put
\begin{equation}\label{eq:value}
  \operatorname{val}(\dsn)\;:=\;\max_{\lvert C\rvert=k}\lmin(\SC),
\end{equation}
a maximum over a finite set.  Since $\SC\psd sI_k$ is equivalent to
$\lmin(\SC)\ge s$, the design $\dsn$ admits a dominating $k$-subset if and
only if $\operatorname{val}(\dsn)\ge1$.
We say $\dsn$ \emph{has value at most} $b$ when $\SC-sI_k$ fails to be
positive semidefinite for every $k$-subset $C$ and every $s>b$; this is
$\operatorname{val}(\dsn)\le b$.  Finally set
\[
  \alpha_k(\FF)\;:=\;\inf_{m}\ \inf_{\dsn}\ \operatorname{val}(\dsn),
\]
the infimum over all sizes and all designs of size $m$ and rank $k$ over
$\FF$.  Thus $\Gtz(m,k)$ holds for every $m$ precisely when
$\alpha_k(\R)\ge1$, and likewise over $\C$; and since
Theorem~\ref{thm:corankone}(iii)--(iv) produces designs of value exactly
$1$ at every rank, the conjecture over $\R$ is the assertion
$\alpha_k(\R)=1$ for every $k$.

At every rank and over either field the value has a floor, which limits
how far the complex counterexamples can push it.

\begin{proposition}[Lean]\label{prop:rankfloor}
Every weighted design of size $m$ and rank $k$ over $\FF$ admits a
$k$-subset $C$ with $\SC\psd\frac1k I_k$.  Hence $\alpha_k(\FF)\ge1/k$.
\end{proposition}

\begin{proof}
Let $A$ be the $m\times k$ matrix whose $c$-th row is $g_c\adj$, so that
$(Ax)_c=\langle g_c,x\rangle$, and let $D=\diag(t_1,\dots,t_m)$.
Condition \eqref{eq:parseval} reads $A\adj DA=I_k$; in particular $A$ has
rank $k$, so some $k$-subset indexes an invertible $k\times k$ submatrix.
Choose $C$ maximizing $\lvert\det A_C\rvert$ among all $k$-subsets, put
$M:=A_C$ and $B:=AM^{-1}$.  Then $A=BM$, so the $c$-th row of $B$ holds
the coefficients expressing $g_c\adj$ in the rows of $M$, and the rows of
$B$ indexed by $C$ form $I_k$.  By Cramer's rule, $B_{ci}\det M$ is the
determinant of the matrix obtained from $M$ by substituting $g_c\adj$ for
its $i$-th row; that matrix is, up to a permutation of rows, $A_{C'}$ with
$C':=(C\setminus\{c_i\})\cup\{c\}$, where $c_i$ is the $i$-th element of
$C$; if $c$ already lies in $C\setminus\{c_i\}$ then that matrix has two
equal rows and its determinant vanishes.  Hence
\[
  \lvert B_{ci}\rvert\;=\;\frac{\lvert\det A_{C'}\rvert}{\lvert\det A_C\rvert}
  \;\le\;1
\]
for every $c$ and $i$, by maximality of the volume \cite{GT2001}.  By
Cauchy--Schwarz, for every $u\in\FF^k$,
\[
  \lvert (Bu)_c\rvert^{2}
  \;\le\;\Bigl(\sum_{i\le k}\lvert B_{ci}\rvert^{2}\Bigr)\lVert u\rVert^{2}
  \;\le\;k\lVert u\rVert^{2} .
\]
Fix $x\in\FF^k$ and take $u:=Mx$, so that $Ax=Bu$ and
$\lVert u\rVert^2=\sum_{c\in C}\lvert\langle g_c,x\rangle\rvert^2$.  Then
\[
  \lVert x\rVert^{2}
  \;=\;x\adj A\adj DAx
  \;=\;\sum_c t_c\lvert (Bu)_c\rvert^{2}
  \;\le\;k\lVert u\rVert^{2}\sum_c t_c
  \;=\;k\sum_{c\in C}\lvert\langle g_c,x\rangle\rvert^{2},
\]
which is \eqref{eq:cover} at $s=1/k$.
\end{proof}

\subsection{The positive half}\label{sec:cpositive}

\begin{proposition}[Lean]\label{prop:cpositive}
Let $k\ge1$ and $m\le k+1$.  Then $\GtzC(m,k)$ holds, as does
$\Gtz(m,k)$.
\end{proposition}

\begin{proof}
The vectors of a design span, so $m\ge k$ and there is nothing to prove
when $m<k$: no design exists.  If $m=k$ the assertion is
Proposition~\ref{prop:sizerank}, and if $m=k+1$ it is
Theorem~\ref{thm:corankone}.  Neither argument uses the field: the first
is the inequality $t_c\le1$ applied to \eqref{eq:parseval}, the second is
the rank-one Schur criterion together with the fact that a weighted mean
does not exceed its maximum.
\end{proof}

By Theorem~\ref{thm:corankone}(iii)--(iv) the bound $1$ is attained at
the last size at which it holds, and attained on a set of full dimension
in the weights: over either field and for every $k\ge1$, every weight
vector $t$ on $k+1$ atoms carries a design at $(k+1,k)$ in which every
$k$-subset dominates and none dominates strictly.  Its chart is
$P=I_{k+1}-zz\adj$ with
$\lvert z_c\rvert^{2}=(1-t_c)/k$, and its leverages obey the tie law
\eqref{eq:tielaw}.

One size lower the value overshoots: for $k\ge2$ every design at $(k,k)$
has value strictly above $1$, and $1$ is approached only as
$\max_c t_c\to1$.

\begin{proposition}\label{prop:square}
Let $\dsn$ be a design at $(k,k)$ over $\FF$.  Then $\ell_c=1/t_c$ for every
$c$, the unique $k$-subset satisfies
\[
  \lmin\bigl(S_{\{1,\dots,k\}}\bigr)\;=\;\frac{1}{\max_c t_c},
\]
and this quantity lies in $(1,k]$ for $k\ge2$.  Every weight vector is
realized, so for $k\ge2$ the infimum of $\operatorname{val}$ over designs
at $(k,k)$ equals $1$ and is not attained.
\end{proposition}

\begin{proof}
In the chart of \S\ref{sec:chart} the matrix $V$ is square with
$V\adj V=I_k$, hence unitary, so its rows are orthonormal:
$\langle v_c,v_d\rangle=\delta_{cd}$ with $v_c=\sqrt{t_c}\,g_c$.  Thus
$\ell_c=\lVert v_c\rVert^2/t_c=1/t_c$ and
\[
  S_{\{1,\dots,k\}}\;=\;\sum_c g_c^{\vphantom{\mathsf{H}}}g_c\adj
  \;=\;\sum_c \frac{1}{t_c}\,v_c^{\vphantom{\mathsf{H}}}v_c\adj,
\]
which is diagonal in the orthonormal basis $(v_c)$ with eigenvalues
$1/t_c$.  Its least eigenvalue is $1/\max_c t_c$.  Since $\sum_c t_c=1$
over $k\ge2$ terms, $1/k\le\max_c t_c<1$, whence the stated range.
Conversely $g_c:=t_c^{-1/2}e_c$ is a design at $(k,k)$ with weights $t$,
and letting $\max_c t_c\to1$ drives the value to $1$ from above.
\end{proof}

\subsection{The seed}\label{sec:cseed}

Let
$\omega:=e^{2\pi i/3}$, so that $1+\omega+\omega^{2}=0$ and
$\omega+\bar\omega=-1$.

\begin{example}[Lean]\label{ex:sic}
Put
\begin{equation}\label{eq:sic}
  g_0:=\bigl(\sqrt2,\,0\bigr),
  \qquad
  g_j:=\Bigl(\sqrt{\tfrac23},\ \sqrt{\tfrac43}\,\omega^{\,j-1}\Bigr)
  \quad(j=1,2,3),
  \qquad t_c:=\tfrac14 .
\end{equation}
Each vector has $\ell_c=2$: for $j\ge1$, $\tfrac23+\tfrac43=2$.  The atom
sum is
\[
  \sum_{j=1}^{3} g_j^{\vphantom{\mathsf{H}}}g_j\adj
  \;=\;\begin{pmatrix}2&\sqrt{8/9}\,\sum_j\overline{\omega^{\,j-1}}\\[2pt]
        \sqrt{8/9}\,\sum_j\omega^{\,j-1}&4\end{pmatrix}
  \;=\;\begin{pmatrix}2&0\\0&4\end{pmatrix},
\]
by $1+\omega+\omega^{2}=0$, so
$\sum_c g_c^{\vphantom{\mathsf{H}}}g_c\adj=4I_2$ and \eqref{eq:parseval} holds at
the uniform weight $\tfrac14$.  This is the $\C^2$ symmetric
informationally complete configuration: four equiangular lines, and indeed
$\lvert\langle g_a,g_b\rangle\rvert^{2}=\tfrac43$ for every $a\ne b$.  For
$b=0$ this is $(\sqrt2\cdot\sqrt{2/3})^2=4/3$; for
$1\le a<b\le3$ it is
\[
  \bigl(\tfrac23+\tfrac43\omega^{\,r}\bigr)
  \bigl(\tfrac23+\tfrac43\overline{\omega^{\,r}}\bigr)
  \;=\;\tfrac49-\tfrac89+\tfrac{16}9\;=\;\tfrac43,
  \qquad r=b-a\not\equiv0,
\]
using $\omega^{r}+\overline{\omega^{r}}=-1$.
\end{example}

\begin{lemma}[Lean]\label{lem:sicvalue}
The design \eqref{eq:sic} has value exactly
\[
  \alpha_{\mathrm{SIC}}\;:=\;2-\frac{2}{\sqrt3}\;=\;0.8452994616\ldots,
\]
the least root of $z^{2}-4z+\tfrac83$.  In particular $\GtzC(4,2)$ is
false, and $\alpha_2(\C)\le\alpha_{\mathrm{SIC}}<\tfrac9{10}$.
\end{lemma}

\begin{proof}
Let $C=\{a,b\}$ be a pair and let $A$ be the $2\times2$ matrix with rows
$g_a\adj,g_b\adj$, so that $\SC=A\adj A$ and
\[
  \det\SC=\lvert\det A\rvert^{2}=\det AA\adj
  =\ell_a\ell_b-\lvert\langle g_a,g_b\rangle\rvert^{2}
  =4-\tfrac43=\tfrac83 ,
  \qquad \tr\SC=\ell_a+\ell_b=4 .
\]
Hence $\det(\SC-zI_2)=z^{2}-4z+\tfrac83$ for every one of the six pairs,
with roots $2\pm2/\sqrt3$, and $\lmin(\SC)=\alpha_{\mathrm{SIC}}$.  The value is the
maximum of these six equal numbers.  Since $\alpha_{\mathrm{SIC}}<1$ no pair dominates.
The inequality $\alpha_{\mathrm{SIC}}<\tfrac9{10}$ is $\tfrac{11}{10}<\tfrac2{\sqrt3}$,
that is $11\sqrt3<20$, that is $363<400$.
\end{proof}

Here the two fields part: over $\R$ rank two holds at every size
(Theorem~\ref{thm:ranktwo}), so $\alpha_2(\R)=1$, while over $\C$ the
same cell has value below $1$ already at size four.  No proof of $\Gtz$
can therefore be field-agnostic.  The constant $\alpha_{\mathrm{SIC}}$
is moreover sharp: by \cite{Nesterenko2026complex} no uniform-weight
design of rank two has value below $2-2/\sqrt3$, the bound being attained exactly when
$4\mid m$, at a repeated copy of \eqref{eq:sic}.  The bound passes to
arbitrary weights in two steps.  Merging equal vectors turns a uniform
design of larger size into one with rational weights, so by
Lemma~\ref{lem:split} read backwards it has the same value and the bound
holds there.  For arbitrary weights, hold the chart $P$ fixed: by
Lemma~\ref{lem:realize} every weight vector $t'$ with positive entries
carries a design with that chart, namely $g'_c=\sqrt{t_c/t'_c}\,g_c$, and
$\operatorname{val}$ is a continuous function of $t'$ there, being a
maximum of finitely many least eigenvalues of matrices depending
continuously on $t'$.  Rational weight vectors are dense in the open
simplex, so the bound persists in the limit.  Hence
$\alpha_2(\C)=2-2/\sqrt3$.

\subsection{Two padding operations}\label{sec:cpadding}

The first operation adjoins one coordinate together with one atom along
it, raising the size and the rank by one.  The second lists an existing
atom twice, raising the size alone.  Each preserves \eqref{eq:parseval}
by a one-line computation, and each preserves a value bound.

\begin{lemma}[Lean]\label{lem:spike}
Let $\dsn$ be a complex weighted design at $(m,k)$ and let $0<w<1$.  Define
$\widehat{\dsn}$ at $(m+1,k+1)$ by
\begin{equation}\label{eq:spike}
  \widehat g_c:=\Bigl(\frac{g_c}{\sqrt{1-w}},\,0\Bigr)\ \ (c\le m),
  \qquad
  \widehat g_{m+1}:=\frac{1}{\sqrt w}\,e_{k+1},
  \qquad
  \widehat t_c:=(1-w)t_c,\quad \widehat t_{m+1}:=w .
\end{equation}
Then $\widehat{\dsn}$ is a design, and if $\dsn$ has value at most $b\ge0$ then
$\widehat{\dsn}$ has value at most $b/(1-w)$.
\end{lemma}

\begin{proof}
The weights are positive and sum to $(1-w)+w=1$.  For the Parseval
condition, the first $m$ terms contribute
\[
  \sum_{c\le m}(1-w)t_c\,
  \frac{1}{1-w}\begin{pmatrix}g_c^{\vphantom{\mathsf{H}}}g_c\adj&0\\0&0\end{pmatrix}
  \;=\;\begin{pmatrix}I_k&0\\0&0\end{pmatrix},
\]
and the last contributes $w\cdot w^{-1}e_{k+1}^{\vphantom{\mathsf{H}}}e_{k+1}\adj$;
the sum is $I_{k+1}$.

Now let $C$ be a $(k+1)$-subset and $s>b/(1-w)$; we show $S_C-sI_{k+1}$ is
not positive semidefinite.  Suppose first $m+1\notin C$.  Every
$\widehat g_c$ with $c\le m$ vanishes in the last coordinate, so the
$(k+1,k+1)$ entry of $S_C-sI_{k+1}$ is $-s<0$, and a positive semidefinite
matrix has nonnegative diagonal.  Suppose instead $m+1\in C$ and put
$C':=C\setminus\{m+1\}$, a $k$-subset of $\{1,\dots,m\}$.  Test on the
flat directions $\widehat x:=(x,0)$ with $x\in\C^k$.  There
\[
  \langle\widehat g_{m+1},\widehat x\rangle=0,
  \qquad
  \langle\widehat g_c,\widehat x\rangle
  =\frac{\langle g_c,x\rangle}{\sqrt{1-w}}\quad(c\le m),
  \qquad
  \lVert\widehat x\rVert=\lVert x\rVert .
\]
If $S_C\psd sI_{k+1}$ then, by the right-hand form of \eqref{eq:cover}
applied to $\widehat x$,
\[
  s\lVert x\rVert^{2}
  \;\le\;\sum_{c\in C}\lvert\langle\widehat g_c,\widehat x\rangle\rvert^{2}
  \;=\;\frac{1}{1-w}\sum_{c\in C'}\lvert\langle g_c,x\rangle\rvert^{2}
  \qquad\text{for every }x\in\C^{k},
\]
that is $S_{C'}\psd (1-w)s\,I_k$ in the original design.  But
$(1-w)s>b$, contradicting the hypothesis on $\dsn$.
\end{proof}

The spike is invisible in the flat directions and is the only atom
visible in the new one.  The factor $1/(1-w)$ costs something real: for
$b>0$ every admissible $w$ has $b<b/(1-w)$, so no member of the family
attains the level $b$, though the levels descend to it as $w\to0$.

\begin{lemma}[Lean]\label{lem:split}
Let $\dsn$ be a complex weighted design at $(m,k)$, let $c_0\le m$ and let
$0<f<1$.  Define $\widetilde{\dsn}$ at $(m+1,k)$ by listing the vector
$g_{c_0}$ twice,
\begin{equation}\label{eq:split}
  \widetilde g_c:=g_c\ (c\le m),
  \quad \widetilde g_{m+1}:=g_{c_0},
  \qquad
  \widetilde t_{c_0}:=f\,t_{c_0},
  \quad \widetilde t_{m+1}:=(1-f)t_{c_0},
  \quad \widetilde t_c:=t_c\ \text{otherwise.}
\end{equation}
Then $\widetilde{\dsn}$ is a design, and if $\dsn$ has value at most $b\ge0$ then
so does $\widetilde{\dsn}$.
\end{lemma}

\begin{proof}
The weights are positive and still sum to one, and in
\eqref{eq:parseval} the term $t_{c_0}g_{c_0}^{\vphantom{\mathsf{H}}}g_{c_0}\adj$ is
replaced by
$f t_{c_0}g_{c_0}^{\vphantom{\mathsf{H}}}g_{c_0}\adj+(1-f)t_{c_0}
g_{c_0}^{\vphantom{\mathsf{H}}}g_{c_0}\adj$, which is the same matrix.

Let $C$ be a $k$-subset of $\{1,\dots,m+1\}$ and $s>b\ge0$.  If $C$ fails
to contain both labels $c_0$ and $m+1$, then replacing $m+1$ by $c_0$ where
necessary carries $C$ to a $k$-subset $C'$ of $\{1,\dots,m\}$ with the
same family of vectors, so $S_C=S_{C'}$ and $S_C-sI_k$ is not positive
semidefinite by the hypothesis on $\dsn$.  If $C$ contains both labels, its
$k$ vectors include two equal ones, hence span a subspace of dimension at
most $k-1$; choosing $x\ne0$ orthogonal to that subspace gives
$x\adj S_Cx=0$ by \eqref{eq:cover}, while
$x\adj(sI_k)x=s\lVert x\rVert^{2}>0$.
\end{proof}

\subsection{The ladder}\label{sec:cladder}

\begin{theorem}[Lean]\label{thm:ladder}
Let $k\ge2$ and $m\ge k+2$.  There is a complex weighted design of size
$m$ and rank $k$ whose value is at most
\[
  \tfrac{10}9\,\alpha_{\mathrm{SIC}}\;=\;\frac{20}{9}-\frac{20\sqrt3}{27}
  \;=\;0.9392216240\ldots\;<\;1 .
\]
In particular $\GtzC(m,k)$ is false.
\end{theorem}

\begin{proof}
Set $e:=k-2\ge0$.  If $e=0$ start with the design \eqref{eq:sic} itself,
of size $4=k+2$ and value $\alpha_{\mathrm{SIC}}$ by Lemma~\ref{lem:sicvalue}.  If
$e\ge1$, apply Lemma~\ref{lem:spike} $e$ times to \eqref{eq:sic} with the
common spike weight
\[
  w\;:=\;1-\Bigl(\tfrac9{10}\Bigr)^{1/e},
  \qquad\text{so that}\qquad (1-w)^{e}=\tfrac9{10}.
\]
Each application raises the size and the rank by one, so after $e$ steps
the size is $4+e=k+2$ and the rank is $2+e=k$; and each application
divides the value bound by $1-w$, so the resulting bound is
\[
  \frac{\alpha_{\mathrm{SIC}}}{(1-w)^{e}}\;=\;\frac{10}{9}\,\alpha_{\mathrm{SIC}} .
\]
This is the asserted bound, and it is below $1$ because
$\alpha_{\mathrm{SIC}}<\tfrac9{10}$ (Lemma~\ref{lem:sicvalue}).  Finally apply
Lemma~\ref{lem:split} $m-k-2$ times, which raises the size to $m$ without
raising the rank and without raising the value bound.  A design of value
below $1$ has no dominating $k$-subset.
\end{proof}

At $k=3$ the theorem is the one-spike padded configuration of size five,
whose value is exactly $\tfrac{10}9\alpha_{\mathrm{SIC}}$; that cell, $(5,3)$, is the
smallest at which rank three fails over $\C$, since
Proposition~\ref{prop:cpositive} covers $m\le4$.  The bound is not the
best available at larger sizes: \S\ref{sec:csharp} exhibits designs of
rank three with value $0.7639\ldots$ and $0.7018\ldots$.

\begin{theorem}[Lean]\label{thm:complex}
Let $k\ge2$.  Then
\[
  \GtzC(m,k)\quad\Longleftrightarrow\quad m\le k+1 .
\]
For $k\le1$, $\GtzC(m,k)$ holds for every $m$.
\end{theorem}

\begin{proof}
For $m\le k+1$ this is Proposition~\ref{prop:cpositive}; for $m\ge k+2$ it
is Theorem~\ref{thm:ladder}.  At $k=1$ the statement is
Proposition~\ref{prop:rankone}, and at $k=0$ the empty subset dominates
the $0\times0$ identity.
\end{proof}

For $k\ge2$ the equivalence is vacuous below $m=k$, where Parseval admits
no design: the assertion holds there for want of anything to
quantify over, not because selection succeeds.

\subsection{The sharp constants}\label{sec:csharp}

Theorem~\ref{thm:complex} settles the decision problem over $\C$ and
leaves the quantitative one.  By Proposition~\ref{prop:rankfloor} and
Theorem~\ref{thm:ladder},
\[
  \frac1k\;\le\;\alpha_k(\C)\;\le\;\frac{10}{9}\alpha_{\mathrm{SIC}}\;<\;1
  \qquad(k\ge2),
\]
and at rank two both ends are known: $\alpha_2(\C)=2-2/\sqrt3$, the upper
bound by Lemma~\ref{lem:sicvalue} and the lower by
\cite{Nesterenko2026complex}.  At rank three the two ends are far apart, and two designs push the
upper end below the padded seed.

\begin{example}[Lean]\label{ex:trine}
Two coplanar trines of lines in $\C^{3}$ sharing an axis: for $j=0,1,2$,
\[
  g_j:=\bigl(\sqrt2,\,0,\,\omega^{\,j}\bigr),
  \qquad
  g_{3+j}:=\bigl(0,\,\sqrt2,\,\omega^{\,j}\bigr),
  \qquad t_c:=\tfrac16 .
\]
Every leverage is $2+1=3$, and
$\sum_{j}g_j^{\vphantom{\mathsf{H}}}g_j\adj=\diag(6,0,3)$,
$\sum_{j}g_{3+j}^{\vphantom{\mathsf{H}}}g_{3+j}\adj=\diag(0,6,3)$, the off-diagonal
entries cancelling by $1+\omega+\omega^{2}=0$; the total is $6I_3$, so
\eqref{eq:parseval} holds.  This design has value exactly $3-\sqrt5$.
\end{example}

To prove the last assertion, and to prepare the rank-three record, we
compute the determinant of a triple once and for all.  Let $C=\{a,b,c\}$ and let $A$ be the
$3\times3$ matrix with rows $g_a\adj,g_b\adj,g_c\adj$.  Then $\SC=A\adj A$
while the Gram matrix is $\Gram_C=AA\adj$, with entries
$\langle g_i,g_j\rangle$; for square $A$ the two have the same
characteristic polynomial.  Writing $\beta_{ij}:=\langle g_i,g_j\rangle$
and expanding the $3\times3$ determinant of
$\Gram_C-sI_3$ gives, with $d:=3-s$ when all three leverages equal $3$,
\begin{equation}\label{eq:tripledet}
  \det(\SC-sI_3)
  \;=\;d^{3}
      -d\bigl(\lvert\beta_{ab}\rvert^{2}+\lvert\beta_{bc}\rvert^{2}
              +\lvert\beta_{ca}\rvert^{2}\bigr)
      +2\operatorname{Re}\beta_{abc},
  \qquad
  \beta_{abc}:=\beta_{ab}\beta_{bc}\beta_{ca}.
\end{equation}
The three moduli form a phase-free budget, and the phases enter through
the single real number $\operatorname{Re}\beta_{abc}$, affinely.

For the trine, the two coplanar triples $\{g_0,g_1,g_2\}$ and
$\{g_3,g_4,g_5\}$ have atom sums $\diag(6,0,3)$ and $\diag(0,6,3)$, of
least eigenvalue $0$.  Every other triple consists of two vectors from one
side and one from the other, say $g_i,g_j,g_{3+l}$ with $i\ne j$; then
\[
  \beta_{ij}=2+\omega^{\,j-i},
  \qquad
  \beta_{j,3+l}=\omega^{\,l-j},
  \qquad
  \beta_{3+l,i}=\omega^{\,i-l},
\]
so the budget is $\lvert2+\omega^{\,j-i}\rvert^{2}+1+1=3+2=5$ and
\[
  \beta_{abc}=(2+\omega^{\,j-i})\,\omega^{\,i-j}=2\omega^{\,i-j}+1,
  \qquad \operatorname{Re}\beta_{abc}=1+2\cdot(-\tfrac12)=0 .
\]
By \eqref{eq:tripledet}, $\det(\SC-sI_3)=d^{3}-5d=d(d^{2}-5)$, so the
spectrum of $\SC$ is $\{3,\,3-\sqrt5,\,3+\sqrt5\}$ and
$\lmin(\SC)=3-\sqrt5=0.7639320225\ldots$ for all eighteen such triples.
The value of the design is the maximum over the twenty triples, namely
$3-\sqrt5$.  In particular $\GtzC(6,3)$ fails, at a level well below the
one Theorem~\ref{thm:ladder} provides.

\begin{example}[Lean]\label{ex:hesse}
The Hesse configuration.  Index nine directions of $\C^{3}$ by a slot
$p\in\{0,1,2\}$ and a phase $q\in\{0,1,2\}$: the direction $u_{3p+q}$
carries $0$ in slot $p$, $\omega^{q}$ in slot $p+1$ and $-\omega^{2q}$ in
slot $p+2$, indices modulo three.  Explicitly, with $\omega$ written $w$,
\[
  \begin{array}{lll}
    u_0=(0,1,-1), & u_3=(-1,0,1), & u_6=(1,-1,0),\\
    u_1=(0,w,-w^{2}), & u_4=(-w^{2},0,w), & u_7=(w,-w^{2},0),\\
    u_2=(0,w^{2},-w), & u_5=(-w,0,w^{2}), & u_8=(w^{2},-w,0).
  \end{array}
\]
Each $u_c$ has $\lVert u_c\rVert^{2}=2$, and
$\sum_c u_c^{\vphantom{\mathsf{H}}}u_c\adj=6I_3$: the family $p$ contributes $1$ to
each of the diagonal slots $p+1$ and $p+2$ for each of its three phases,
so every slot receives $6$ from the two families that do not vanish there,
while its one off-diagonal entry,
$\omega^{q}\overline{(-\omega^{2q})}=-\omega^{-q}$ in position
$(p+1,p+2)$, sums to $-\sum_q\omega^{-q}=0$ over the family.  Put
\[
  g_c:=\sqrt{\tfrac32}\,u_c,\qquad t_c:=\tfrac19 .
\]
Then \eqref{eq:parseval} holds, $\ell_c=3$ for every $c$, and for $a\ne b$
\[
  \langle u_a,u_b\rangle^{3}=-1,
  \qquad\text{hence}\qquad
  \beta_{ab}^{3}=\langle g_a,g_b\rangle^{3}=-\tfrac{27}8 .
\]
The overlap identity is two computations.  Within the family $p$, with
$r:=q'-q\not\equiv0$,
\[
  \langle u_{3p+q},u_{3p+q'}\rangle=\omega^{r}+\omega^{2r}=-1 ;
\]
and between the families $p$ and $p+1$ the only slot on which both are
nonzero is $p+2$, so
\[
  \langle u_{3p+q},u_{3(p+1)+q'}\rangle
  =\overline{(-\omega^{2q})}\,\omega^{q'}=-\omega^{\,q+q'} .
\]
Each value is $-1$ times a cube root of unity, hence a cube root of $-1$;
the remaining pairs are the conjugates of these.
\end{example}

Since $\lvert\beta_{ab}\rvert^{2}$ is a nonnegative real whose cube is
$\lvert\beta_{ab}^{3}\rvert^{2}=(27/8)^{2}$, it equals $9/4$; the
configuration is equiangular, and the budget in \eqref{eq:tripledet} is
$3\cdot\tfrac94=\tfrac{27}4$ for every triple.  The triangle invariant
satisfies $\beta_{abc}^{3}=(\beta_{ab}\beta_{bc}\beta_{ca})^{3}=-(27/8)^{3}$, so $\beta_{abc}$
is one of the three cube roots of $-(27/8)^3$, namely $-\tfrac{27}8$ and
$\tfrac{27}{16}(1\pm i\sqrt3)$.  Hence
\begin{equation}\label{eq:hessedet}
  \det(\SC-sI_3)\;=\;(3-s)^{3}-\tfrac{27}4(3-s)+2\operatorname{Re}\beta_{abc},
  \qquad \operatorname{Re}\beta_{abc}\in\Bigl\{-\tfrac{27}8,\ \tfrac{27}{16}\Bigr\},
\end{equation}
the first alternative occurring exactly when $\det\SC=\tfrac{27}4+2\operatorname{Re}\beta_{abc}$
vanishes, i.e.\ for the twelve triples lying on a line of the Hesse
configuration, and the second for the remaining seventy-two.  With
$\operatorname{Re}\beta_{abc}=\tfrac{27}{16}$ the right-hand side of
\eqref{eq:hessedet} expands to
\begin{equation}\label{eq:hessecubic}
  \det(\SC-sI_3)\;=\;-\tfrac18\,h(s),
  \qquad
  h(s):=8s^{3}-72s^{2}+162s-81 ,
\end{equation}
and with $\operatorname{Re}\beta_{abc}=-\tfrac{27}8$ it is $-s(s-\tfrac92)^{2}$.

\begin{proposition}[Lean]\label{prop:hesse}
The roots of $h$ are $3(1-\cos\tfrac{2\pi}9)$,
$3(1-\cos\tfrac{4\pi}9)$ and $3(1-\cos\tfrac{8\pi}9)$.  The design of
Example~\ref{ex:hesse} has value exactly
\[
  h_3\;:=\;3\Bigl(1-\cos\tfrac{2\pi}9\Bigr)\;=\;0.7018666706\ldots
\]
Consequently $\alpha_3(\C)\le h_3$, and $\GtzC(m,3)$ fails for every
$m\ge5$.
\end{proposition}

\begin{proof}
Substituting $s=3(1-\gamma)$ in \eqref{eq:hessecubic} gives
\[
  h\bigl(3(1-\gamma)\bigr)\;=\;-27\bigl(8\gamma^{3}-6\gamma+1\bigr),
\]
and for $\gamma=\cos\theta$ the identity $4\cos^{3}\theta-3\cos\theta
=\cos3\theta$ turns $8\gamma^{3}-6\gamma+1=0$ into
$\cos3\theta=-\tfrac12$.  As $\theta$ runs over $[0,\pi]$ the angle
$3\theta$ runs over $[0,3\pi]$, where the cosine takes the value
$-\tfrac12$ at $\tfrac{2\pi}3$, $\tfrac{4\pi}3$ and $\tfrac{8\pi}3$; the
solutions are therefore
$\theta=\tfrac{2\pi}9,\tfrac{4\pi}9,\tfrac{8\pi}9$, which gives the three
roots.  They are distinct, and $s=3(1-\cos\theta)$ increases with $\theta$,
so $h_3$ is the least.  Exact evaluation isolates it: $h$ is negative at
$\tfrac7{10}$, where its value is $-\tfrac{136}{1000}$, and positive at
$\tfrac{71}{100}$, where its value is $\tfrac{588088}{10^{6}}$, so
\begin{equation}\label{eq:hessewindow}
  \tfrac7{10}\;<\;h_3\;<\;\tfrac{71}{100}.
\end{equation}

By \eqref{eq:hessecubic} the seventy-two nondegenerate triples all have
characteristic polynomial $-\tfrac18h$, hence
$\lmin(\SC)=h_3$ for each; the twelve degenerate triples have
$\lmin(\SC)=0$.  The value is the maximum, $h_3$.  By
\eqref{eq:hessewindow}, $h_3<1$, so no triple
dominates, so $\GtzC(9,3)$ fails; Lemma~\ref{lem:split} propagates the
failure to every $m\ge9$, and Theorem~\ref{thm:ladder} covers
$5\le m\le8$.
\end{proof}

An explicit certificate can replace the eigenvalue computation,
exhibiting positive semidefiniteness at the threshold without reference
to the spectral theorem.  Take the triple $\{u_0,u_1,u_3\}$.
From $1+\omega+\omega^{2}=0$,
\[
  u_0^{\vphantom{\mathsf{H}}}u_0\adj+u_1^{\vphantom{\mathsf{H}}}u_1\adj
  +u_3^{\vphantom{\mathsf{H}}}u_3\adj
  \;=\;\begin{pmatrix}1&0&-1\\0&2&\omega\\-1&\omega^{2}&3\end{pmatrix},
\]
so that $\SC$ is $\tfrac32$ times this matrix.  Writing
$p:=\tfrac32-s$ and $q:=3-s$ and $\operatorname{atom}(v):=vv\adj$, a
direct expansion of both sides gives the polynomial identity
\begin{equation}\label{eq:hesseldl}
  pq\,\bigl(\SC-sI_3\bigr)
  \;=\;q\cdot\operatorname{atom}\bigl(p,0,-\tfrac32\bigr)
      +p\cdot\operatorname{atom}\bigl(0,q,\tfrac32\omega^{2}\bigr)
      -\tfrac18h(s)\,E_{33},
\end{equation}
$E_{33}$ being the matrix unit in the last diagonal slot.  The identity is
the $LDL\adj$ factorization of $\SC-sI_3$ cleared of denominators: the
leading $2\times2$ block of $\SC-sI_3$ is the diagonal $\diag(p,q)$, the
two rank-one terms are the pivot contributions of that factorization
multiplied by $pq$, and the coefficient of $E_{33}$ is $pq$ times the
Schur complement of the block, which is where $h$ enters.  Checking the
identity is a computation: the entries other than $(3,3)$ agree as
polynomials in $s$, and the residue at $(3,3)$ is
\[
  pq\bigl(\tfrac92-s\bigr)-\tfrac94(p+q)
  \;=\;\Bigl(-s^{3}+9s^{2}-\tfrac{99}4s+\tfrac{81}4\Bigr)
        -\Bigl(\tfrac{81}8-\tfrac92 s\Bigr)
  \;=\;-\tfrac18h(s).
\]
At $s=h_3$ the cubic vanishes and both $p$ and $q$ are positive, since
$h_3<\tfrac{71}{100}<\tfrac32$; dividing \eqref{eq:hesseldl} by the
positive scalar $pq$ exhibits $\SC-h_3I_3$ as a nonnegative combination of
two rank-one positive semidefinite matrices.  The threshold $h_3$ is
therefore attained by this triple.

The three complex records at rank three are ordered
\[
  h_3=0.7018666706\ldots
  \;<\;3-\sqrt5=0.7639320225\ldots
  \;<\;\tfrac{10}9\alpha_{\mathrm{SIC}}=0.9392216240\ldots ,
\]
both inequalities being exact.  For the second, $12<7\sqrt3$ gives
$\alpha_{\mathrm{SIC}}>\tfrac56$ and hence $\tfrac{10}9\alpha_{\mathrm{SIC}}>\tfrac{25}{27}$, while
$5>\bigl(\tfrac{223}{100}\bigr)^{2}$ gives $3-\sqrt5<\tfrac{77}{100}$, and
$2500>2079$.  For the first,
\[
  h\bigl(3-\sqrt5\bigr)\;=\;\bigl(576-1008+486-81\bigr)
   +\bigl(-256+432-162\bigr)\sqrt5\;=\;14\sqrt5-27\;>\;0,
\]
since $980>729$.  The three roots of $h$ sum to $9$, so they cannot all
lie below $2$; with $h(2)=19>0$ this places $2$ strictly between the two
smallest roots.  As $3-\sqrt5<2$ and $h(3-\sqrt5)>0$, the number
$3-\sqrt5$ lies in the same interval, hence above $h_3$.

\begin{remark}\label{rem:phase}
Identity \eqref{eq:tripledet} shows where the two fields differ.
Over $\R$ the entries $\beta_{ij}$ are real, so $\beta_{abc}$ is real and
$\operatorname{Re}\beta_{abc}=\pm\lvert\beta_{abc}\rvert$ is forced to an extreme of its
range; over $\C$ it may sit anywhere in
$[-\lvert\beta_{abc}\rvert,\lvert\beta_{abc}\rvert]$, and the separating
quantity is the phase defect
$\lvert\beta_{abc}\rvert^{2}-(\operatorname{Re}\beta_{abc})^{2}=(\operatorname{Im}\beta_{abc})^{2}$,
which vanishes identically for every real configuration.  The trine of
Example~\ref{ex:trine} isolates the mechanism: its
modulus budget $5$ would permit $2\operatorname{Re}\beta_{abc}=2\sqrt3$ and hence a
positive determinant at $s=1$, but its phase defect is $3$ and its
$\operatorname{Re}\beta_{abc}$ is $0$, leaving $\det(\SC-I_3)=-2$.  No
argument that sees only the moduli $\lvert\beta_{ij}\rvert$ can decide
either problem; the corresponding no-go statement is recorded in
Appendix~\ref{app:graveyard}.
\end{remark}

\begin{remark}\label{rem:sharpc}
Collecting the two ends,
\[
  \tfrac13\;\le\;\alpha_3(\C)\;\le\;3\bigl(1-\cos\tfrac{2\pi}{9}\bigr)
  \;=\;0.7018\ldots,
\]
the lower bound from Proposition~\ref{prop:rankfloor} and the upper from
Proposition~\ref{prop:hesse}.  The gap stays open: every rank-three
statement above bounds $\alpha_3(\C)$ from above, and such bounds move
the record down without pinning it.  Determining $\alpha_3(\C)$ is
Problem~\ref{prob:sharpc}.  The rank-inverse bound $1/k$ is sharp only at
$k=1$, where the one-atom design at $(1,1)$ has value $1$.
\end{remark}
